\BeleskeID{07-s014}
% element: 07-s014:d01
% element: 07-s014:e01
% element: 07-s014:e02
Kao što jezik visokog nivoa izražava složen koncept u programu, VHDL omogućava da se ponašanje složenog elektronskog kola unese u sistem za automatsku sintezu ili simulaciju. Poput Pascal-a, C-a i C++-a, ima sredstva za strukturiran dizajn, upravljanje tokom i predstavljanje podataka. Dodatno mora opisivati istovremene događaje, jer i sam hardver radi istovremeno. Korisnicima PALASM-a, ABEL-a i CUPL-a taj pristup je poznat; studentima koji imaju iskustva samo sa sekvencijalnim softverskim programima to je novi koncept: konkurentni delovi opisa predstavljaju hardver koji deluje istovremeno.

% element: 07-s014:e03
Testbench opisuje pobude kola i odgovarajuće očekivane izlaze tokom vremena, odnosno pretvara zahteve performansi u proveru. Ta mogućnost je značajna i često nedovoljno korišćena. Testbench treba praviti paralelno sa opisom kola, kao deo svakog VHDL projekta, umesto da se provera ostavi tek za kraj projektovanja.
